\centering
\section*{\underline{Ähnlichkeitstransformation}}

\begin{tikzpicture}[remember picture, overlay]

% ---- ERSTE REIHE: Am Seitenursprung starten
\setlength{\rowheight}{40mm}

    % A1) Erste Kachel oben links
    \setlength{\cellwidth}{43mm}
    \Tile{A1}{([xshift=10mm,yshift=-18mm]current page.north west)}{\cellwidth}{\rowheight}
        {
            \vspace{-\TileSep}
            \titlerm{Bedingungen}\\ \vspace{-1mm}
            \begin{equation*}
                \text{Ma} \stackrel{!}{>} \text{0,3}
            \end{equation*}
            {\scriptsize Keine Gültigkeit für den transsonischen und hypersonischen Bereich\\
            bzw. ab Ma $>$ 0,8}
        }{0}{1}{1}{0}

    % B1) Zweite Kachel rechts anbauen
    \setlength{\cellwidth}{147mm}
    \Tile{B1}{(A1.north east)}{\cellwidth}{\rowheight}
        {
            \vspace{-\TileSep}
            \titlebf{Göthert-Regel (3D)}\\ \vspace{2mm}
            \begin{minipage}{0.4\cellwidth-\TileSep}
                \centering
                \titlerm{Allgemein}\\ \vspace{-3mm}
                \begin{equation*}
                    \beta =
                    \begin{cases}
                        \sqrt{1 - \text{Ma}^2_\infty}     & \text{für Ma} < 1 \\
                        \sqrt{\text{Ma}^2_\infty - 1} = B & \text{für Ma} > 1
                    \end{cases}
                \end{equation*}%
                \vspace{-2mm}\\
                {\scriptsize Vergleichsströmung\dots \\[0.8ex]
                \dots \textbf{subsonisch} $ik$ bei Ma$_\infty = 0$ \\
                \dots \textbf{supersonisch} $ik$ bei Ma$_\infty = \sqrt{2}$ \\}
                \vspace{2mm}
            \end{minipage}%
            \begin{minipage}{0.6\cellwidth-\TileSep}
                \centering
                \titlerm{Transformation}\\ \vspace{-5.5mm}
                \begin{align*}
                    \left\{x\right\}_{ik} &\Leftrightarrow \left\{x\right\} \\
                    \left\{y, z, f, \alpha, \delta, \lambda\right\}_{ik} &\Leftrightarrow \beta\cdot\left\{y, z, f, \alpha, \delta, \lambda\right\} \\
                    \left\{C_\text{A}, C_\text{M}, c_p\right\}_{ik} &\Leftrightarrow \beta^2\cdot\left\{C_\text{A}, C_\text{M}, c_p\right\} \\[0.8ex]
                    \left\{x,y,z,\alpha,f,\delta,\Delta\right\}_{ik} &\Leftrightarrow {\underbrace{\tfrac{1}{\beta^2}}_{\mathclap{\text{Rück-Trafo}}}}\,c_{p,ik}\,\{x,{\underbrace{\beta y,\beta z,\beta\alpha,\beta f,\beta\delta,\beta\Delta}_\text{Hin-Trafo}}\}
                \end{align*}
            \end{minipage}
        }{0}{0}{1}{0}

% ---- ZWEITE REIHE: Unten anbauen
\setlength{\rowheight}{39mm}

    % A2) Erste Kachel unten anbauen
    \setlength{\cellwidth}{95mm}
    \Tile{A2}{(A1.south west)}{\cellwidth}{\rowheight}
        {
            \titlebf{Prandtl-Glauert-Regel (2D)}\\ \vspace{1mm}
            {\color{Purple}\titlerm{Unveränderte Geometrie, gut für die Nachrechnung}}\\ \vspace{-1mm}
            \begin{equation*}
                \left\{C_\text{A}, C_\text{M}, c_p\right\}_{ik} \Leftrightarrow \beta\cdot\left\{C_\text{A}, C_\text{M}, c_p\right\}
            \end{equation*}%
            Gesamt-Trafo:\\
            \vspace{-5mm}
            \begin{equation*}
                c_p\!\left(x, z, \alpha, f, \delta\right) = \frac{1}{\beta}\,c_{p,ik}\!\left(x, z, \alpha, f, \delta\right)
            \end{equation*}
        }{0}{1}{1}{0}

    % B2) Zweite Kachel rechts anbauen
    \setlength{\cellwidth}{95mm}
    \Tile{B2}{(A2.north east)}{\cellwidth}{\rowheight}
        {
            \titlebf{Prandtl-Glauert-Regel (2D)}\\ \vspace{1mm}
            {\color{NavyBlue}\titlerm{Unveränderte Druckverteilung, gut für den Entwurf}}\\ \vspace{-2mm}
            \begin{equation*}
                \left\{z, f, \alpha, \delta\right\}_{ik} \Leftrightarrow \frac{1}{\beta}\cdot\left\{z, f, \alpha, \delta\right\}
            \end{equation*}%
            Gesamt-Trafo:\\
            \vspace{-5mm}
            \begin{equation*}
                c_p\!\left(x, z, \alpha, f, \delta\right) = c_{p,ik}\!\left(x, \frac{1}{\beta}z, \frac{1}{\beta}\alpha, \frac{1}{\beta}f, \frac{1}{\beta}\delta\right)
            \end{equation*}
        }{0}{0}{1}{0}

\end{tikzpicture}

\centering
\vspace{74mm}
\section*{\underline{Profilumströmung Überschall}}

\begin{tikzpicture}[remember picture, overlay]

% ---- ERSTE REIHE: Am Seitenursprung starten
\setlength{\rowheight}{50mm}

    % A1) Erste Kachel oben links
    \setlength{\cellwidth}{74mm}
    \Tile{A1}{([xshift=10mm,yshift=-111mm]current page.north west)}{\cellwidth}{\rowheight}
        {
            \vspace{-\TileSep}
            \titlebf{Annahmen}\\ \vspace{1mm}
            \titlerm{Wegen Linearisierung:}\\ \vspace{-1mm}
            \begin{itemize}[leftmargin=5mm]
                \item Alle Stöße sind isentrope Verdichtungs-\\ wellen
                \item Alle Charakteristiken haben den gleichen\\ Winkel ($\mu_\infty$)
                \item Keine Nachlaufdelle, keine Drosselung
            \end{itemize}
            \begin{align*}
                \text{\titlerm{Skelettlinie:}}\quad &\bar{z}_\text{o} = \bar{z}_\text{u} \\
                \text{\titlerm{Profiltropfen:}}\quad &\alpha^* = 0\,;~\bar{z} = \tfrac{1}{2}d(\bar{x})\,;~\bar{z}_\text{o} = -\bar{z}_\text{u}
            \end{align*}
            \\ \vspace{-1mm}
        }{0}{1}{1}{0}

    % B1) Zweite Kachel rechts anbauen
    \setlength{\cellwidth}{116mm}
    \Tile{B1}{(A1.north east)}{\cellwidth}{\rowheight}
        {
            \vspace{-\TileSep}
            \titlerm{Beiwerte am Profil}\\ \vspace{-3mm}
            \begin{align*}
                C_\text{A} &= \frac{4}{\sqrt{\text{Ma}^2_\infty - 1}}\,\int_0^1 \alpha\,\text{d}\bar{x} = \frac{4}{\sqrt{\text{Ma}^2_\infty - 1}}\,\int_0^1 \alpha^* - \frac{\text{d}\bar{z}}{\text{d}\bar{x}}\,\text{d}\bar{x} \\
                C_\text{M} &= -\frac{4}{\sqrt{\text{Ma}^2_\infty - 1}}\,\left(\frac{\alpha^*}{2} - \int_0^1 \frac{\text{d}\bar{z}(\bar{x})}{\text{d}\bar{x}}\,\bar{x}\,\text{d}\bar{x}\right) \\
                C_\text{W,W} &= \frac{4}{\sqrt{\text{Ma}^2_\infty - 1}}\,\left({\alpha^*}^2 + \frac{1}{2}\,\int_0^1\left(\frac{\text{d}\bar{z}_\text{o}}{\text{d}\bar{x}}\right)^2 + \left(\frac{\text{d}\bar{z}_\text{u}}{\text{d}\bar{x}}\right)^2\,\text{d}\bar{x}\right)
            \end{align*}
            \begin{equation*}
                c_p = \frac{4}{B}\left(\alpha^* - \frac{\partial\bar{z}}{\partial\bar{x}}\right)\,; \qquad \frac{\text{d}C_\text{A}}{\text{d}\alpha} = \frac{4}{B} = \text{const.}\,; \qquad \varepsilon = \frac{W_\text{W}}{A} = \frac{C_\text{W,W}}{C_\text{A}} = \alpha
            \end{equation*}
        }{0}{0}{1}{0}

% ---- ZWEITE REIHE: Unten anbauen
\setlength{\rowheight}{34mm}

    % A2) Erste Kachel unten anbauen
    \setlength{\cellwidth}{90mm}
    \Tile{A2}{(A1.south west)}{\cellwidth}{\rowheight}
        {
            \titlerm{Beiwerte der ebenen Platte}\\ \vspace{2mm}
            \begin{minipage}{0.5\cellwidth-\TileSep}
                \begin{align*}
                    C_\text{A} &= \frac{4}{\sqrt{\text{Ma}_\infty^2 - 1}} \\[0.5ex]
                    C_\text{W,W} &= \frac{4\,\alpha^2}{\sqrt{\text{Ma}_\infty^2 - 1}}
                \end{align*}
            \end{minipage}%
            \begin{minipage}{0.5\cellwidth-\TileSep}
                \begin{align*}
                    C_\text{M} &= -\frac{C_\text{A}}{2} + C_\text{M,A=0} \\[0.5ex]
                    c_p &= \frac{\mp2\,\alpha}{\sqrt{\text{Ma}_\infty^2 - 1}}
                \end{align*}
            \end{minipage}
        }{0}{1}{1}{0}

    % B2) Zweite Kachel rechts anbauen
    \setlength{\cellwidth}{100mm}
    \Tile{B2}{(A2.north east)}{\cellwidth}{\rowheight}
        {
            \titlebf{Druckpunkt / Neutralpunkt}\\ \vspace{2mm}
            \begin{minipage}{0.5\cellwidth-\TileSep}
                \centering
                \titlerm{Relative Druckpunktlage}
                \begin{align*}
                    \frac{x_\text{A}}{t} &= -\frac{C_\text{M}}{C_\text{A}} \\[0.5ex]
                                         &= \frac{1}{2} - \frac{C_\text{M,A=0}}{C_\text{A}}
                \end{align*}
            \end{minipage}%
            \begin{minipage}{0.5\cellwidth-\TileSep}
                \centering
                \titlerm{Relative Neutralpunktlage}
                \begin{align*}
                    \frac{x_\text{N}}{t} &= -\frac{\text{d}C_\text{M,VK}}{\text{d}C_\text{A}} \\[0.5ex]
                                         &= \frac{1}{2}
                \end{align*}
            \end{minipage}
        }{0}{0}{1}{0}

\end{tikzpicture}